% titel: Blatt 2, A5: Bäume

% Teil I: Werkzeuge, kein Lösungsweg.

Für einen Baum, in dem jeder Knoten, der kein freies Ende ist, genau drei
benachbarte Kanten hat, zeigen: $n$ ist gerade und der Baum hat genau
$\tfrac n2+1$ freie Enden.

\begin{werkzeugkasten}[Begriffe und Handschlaglemma]
Ein \emph{Baum} ist ein zusammenhängender kreisfreier Graph. Ein
\emph{freies Ende} ist ein Knoten vom Grad $1$ (Blatt); die übrigen Knoten
haben Grad $3$. In jedem Baum gilt \emph{Eckenbilanz}:
\begin{equation}
  \sum_{\text{Knoten }v}\deg(v)=2\cdot(\text{Anzahl Kanten})=2(n-1).
\end{equation}
Diese Bilanz liefert bereits $n$ gerade und $\tfrac n2+1$ freie Enden.
\end{werkzeugkasten}

\begin{werkzeugkasten}[Induktionsschritt mit freien Enden]
Per Induktion über $n$. Für $n>2$ gibt es einen Knoten, der zu
\emph{mindestens} zwei freien Enden benachbart ist. Entfernt man diese
beiden freien Enden zusammen mit den angrenzenden Kanten, wird der nun
freie Knoten selbst ein freies Ende: $n$ sinkt um $2$, die Zahl der freien
Enden um $1$. Damit ist die Zahl $\tfrac n2+1$ invariante:
$(n-2)/2+1=n/2$.
\end{werkzeugkasten}

\begin{achtung}[Darauf achten]
\begin{itemize}
  \item Grad-Zusammenhang: zweimal Kanten zählen (jede Kante hat zwei
        Enden) -- daraus das ``gerade $n$''-Argument ableiten.
  \item Beim induktiven Entfernen \emph{beide} freien Enden samt des
        Zwischenknotens nehmen, sonst zerstört man die ``nur Grad 3''-
        Bedingung der inneren Knoten.
  \item Nicht nur $n$ gerade zeigen, sondern auch die Anzahl der freien
        Enden ($\tfrac n2+1$) -- beides gehört zur Behauptung.
\end{itemize}
\end{achtung}
